\documentclass{article}
% Shared layout. Keep benchmark prose and result-specific plots outside this folder.
\PassOptionsToPackage{numbers}{natbib}
\usepackage[preprint]{neurips_2026}
\usepackage[T1]{fontenc}
\usepackage[hidelinks]{hyperref}
\usepackage{booktabs,amsmath,microtype,pgfplots}
\pgfplotsset{compat=1.17}
\usetikzlibrary{patterns}
\definecolor{okblue}{HTML}{0072B2}
\definecolor{okorange}{HTML}{D55E00}
\pgfplotsset{every axis/.append style={font=\small, tick label style={font=\small}, width=\linewidth, height=5cm, grid=major}}

\newcommand{\TaskCount}{60}
\newcommand{\FamilyCount}{12}
\newcommand{\KindCount}{5}
\newcommand{\VisibleRate}{100.0\% [94.0, 100.0]}
\newcommand{\FloorRate}{0.0\% [0.0, 6.0]}
\newcommand{\StaleRate}{100.0\% [94.0, 100.0]}
\newcommand{\PairDifference}{100.0 [100.0, 100.0]}
\newcommand{\AuthorOrcid}{0009-0005-9964-6841}
\newcommand{\SimAccuracy}{52.0\% [50.8, 53.2]}
\newcommand{\SimBlocked}{48.0\% [46.8, 49.2]}
\newcommand{\SimStaleBefore}{48.0\% [46.8, 49.2]}
\newcommand{\SimStaleAfter}{0.0\% [0.0, 0.0]}
\newcommand{\SimDifference}{0.0 [0.0, 0.0]}
\newcommand{\AdherenceGrid}{0, 0.25, 0.5, 0.75, 1}
\newcommand{\BootstrapDraws}{2000}
\newcommand{\AgentCount}{100}
\newcommand{\SimulationCalls}{60000}
\newcommand{\BaselineCalls}{960}

\title{Schema Mutation: Measuring Tool Calls After Interface Changes}
\author{Pooria Arab\\Independent researcher\\ORCID: \AuthorOrcid}
\begin{document}
\maketitle
\begin{abstract}Tool interfaces change while familiar call patterns remain. Schema Mutation measures whether a call follows the current interface after a migration. The benchmark crosses 12 tool families with 5 mutation types, yielding 60 public tasks. A prompt-only scripted solver achieves 100.0\% accuracy (95\% Wilson interval 94.0-100.0\%). Empty, constant, refusal, echo, and first-tool controls score 0.0\% (0.0-6.0\%). Exact stale calls fail and receive a separate diagnostic. A simulated population sweep compares pass-through with schema validation. At adherence q=0.5, simulated accuracy is 52.0\% (agent-bootstrap interval 50.8-53.2\%) under both layers. Validation blocks stale proposals without improving completion. These are scripted and simulated results, not learned-model evaluations. The study supports measurement checks on a synthetic public task set; it does not establish deployment reliability or identify memorization as a cause.
\end{abstract}
\section{Why interface changes matter}
An agent can understand a request and still issue a call that no longer exists. A field may change its name. A value may acquire a different type. An endpoint may split into separate operations. Familiar examples can then become a source of error. The current interface is available, but the agent must apply it to the request before acting.\par
Schema Mutation measures current-interface adherence after an explicit migration. Its unit of evaluation is a single proposed tool call. The task supplies the current schemas and a user request. The answer must select the correct operation and express the requested values under those schemas. A separate diagnostic identifies an exact match to the known pre-migration call. This diagnostic distinguishes a specific obsolete answer from other failures. It does not identify the mental process that caused the answer.\par
The scientific claim is deliberately narrow. The benchmark offers controlled, synthetic instances of interface change with deterministic scoring. Scripted controls test whether the measurement separates current calls, obsolete calls, and several unproductive responses. An additional simulated experiment tests whether completion and stale-call suppression respond differently to an enforcement layer. Neither experiment measures a trained language model. Neither supports a ranking of commercial systems.\par
The motivating boundary is between semantic intent and interface representation. A request to refund money can remain unchanged while the amount moves from dollars to cents. A request to send a message can remain unchanged while the recipient field changes its name. In both cases, preserving the request requires changing the call. Measuring fluent text alone misses this boundary. Measuring only whether the answer parses also misses it.\par
A useful evaluation must therefore report successful completion and a specific error signal together. An agent that emits no call cannot use an obsolete schema, but it also cannot complete the request. An enforcement layer can have the same effect. This paper treats those outcomes separately throughout. The simulated experiment makes the distinction visible under explicit assumptions, rather than attributing it to an unobserved learned policy.\par
\section{Related work and measurement scope}
Function-calling evaluations provide the closest comparison. The Berkeley Function Calling Leaderboard studies function selection and argument construction across varied settings \citep{bfcl}. Schema Mutation narrows that broader target to interface changes. Each item retains an explicit obsolete call as a diagnostic reference. The resulting distinction is between a correct current call, an exact old call, and another answer. This paper does not claim that other function-calling suites contain no interface variation.\par
Task completion in a longer interaction requires more than a correct call. The $\tau$-bench evaluation uses tool-mediated conversations, domain policies, and final database states \citep{tau}. Schema Mutation removes those interaction demands to isolate a local conversion problem. It does not test persistence across turns, customer clarification, state recovery, or the consequences of execution. The narrower scope makes failures easier to inspect but limits external validity.\par
Benchmark quality frameworks motivate the control battery. BetterBench emphasizes the design, implementation, documentation, and maintenance of an evaluation \citep{betterbench}. The Agentic Benchmark Checklist highlights task setup and reward validity \citep{abc}. Here, an oracle response checks the positive scoring path. Several unproductive controls check negative paths. A visible-information solver adds evidence that the expected answer is recoverable from the supplied task, without handing that solver the key.\par
PayBench provides a useful structural contrast \citep{paybench}. Its public description distinguishes acting when permitted from stopping when a scenario contains a constraint. Schema Mutation instead asks for a compliant call on every item. Its obsolete call is a negative reference for the same request, not a separately administered harmless scenario. Consequently, this paper uses current-call accuracy and stale-call frequency. It does not import a matched-pair payment metric or a human-preference score.\par
These comparisons suggest a specific use. Schema Mutation can accompany a broader evaluation when interface adherence is a suspected failure source. It cannot replace an interactive task benchmark. Good results here establish compatibility with a small synthetic contract. They do not establish that an agent chooses worthwhile actions, obtains permission, or verifies the effects of its tools.\par
\section{Experiment design}
\subsection{Construct and unit of analysis}
The construct is adherence to the current tool interface when a migration makes an older call incorrect. Each task contains a current schema, a request, a canonical current call, and a canonical obsolete call. Only the schema and request belong in the solver input. The answer references support scoring and diagnostic controls. This separation matters because a solver that reads the answer key can validate the grader but cannot establish task solvability from visible information.\par
The public set contains \TaskCount{} tasks. It crosses \FamilyCount{} tool families with \KindCount{} migration types. Each family contributes an item for each type. A task is a single-turn call-construction problem. The benchmark does not execute a real service or observe an external state transition. The measured outcome is compatibility with the specified call, after documented normalization.\par
The same task supports a positive and a negative reference. The positive reference expresses the request under the current schema. The negative reference expresses it under the old schema. These are paired answers, not paired environments. A current-schema result therefore cannot reveal the size of a model's migration penalty relative to an unmigrated condition. Measuring that penalty would require an additional, matched baseline task with the old interface still valid.\par
\subsection{Scripted controls}
The visible solver reads the prompt and uses a fixed parser for the authored request patterns. It matches values to property descriptions, applies explicit enum maps, converts supported types, and chooses among split endpoints. It does not read the expected call or obsolete-call metadata. This is a specialized scripted solver, not a general language-understanding system. Its success supports recoverability within the authored task grammar.\par
The oracle returns the expected call directly. The stale control returns the obsolete call directly. Both intentionally use privileged information. Their purpose is to exercise the scoring boundary. The oracle asks whether an intended answer receives credit. The stale control asks whether the obsolete reference fails and receives the specific diagnostic. Treating either as a model baseline would misstate what was tested.\par
The remaining controls return an empty answer, a fixed no-operation call, a refusal, the prompt itself, or the first listed tool with empty arguments. They probe inactivity, irrelevant structure, non-completion, copying, and incomplete tool selection. These controls are useful because they can expose credit unrelated to task completion. Passing them is only a limited check: the battery does not enumerate every possible exploit or malformed response.\par
Each control traverses the public benchmark runner twice with the same task set and seed. The study retains final answers, scores, and diagnostic fields. Repeat checks compare semantic records after removing timing information. Repetition checks reproducibility; it does not double the number of independent tasks. The battery contains \BaselineCalls{} scored task executions including repeats.\par
\subsection{Simulated population and control layer}
The extra experiment is explicitly simulated. A simulated agent has an adherence parameter $q$. For each task, a deterministic pseudorandom draw selects the current call with probability parameter $q$, otherwise the obsolete call. The simulated population contains \AgentCount{} seeded agents. Shared draws pair the adherence settings and the control layers. The simulated sweep uses $q \in \{\AdherenceGrid\}$, with integer agent seeds from unity through \AgentCount{}.\par
The simulated pass-through layer emits the proposed call unchanged. The simulated validation layer checks it against the visible schema. It checks the operation name, required fields, extra fields, primitive types, array item types, and enumerated values. If the proposal fails these checks, the layer emits an empty answer. It does not repair the call, consult the expected answer, or retry. This is a shape-validation control, not a general semantic verifier.\par
The simulated experiment deliberately separates proposal behavior from emitted behavior. A blocked obsolete proposal remains evidence of the simulated agent's failure to adapt. The emitted answer, however, is empty and no longer an exact stale call. Reporting only emitted stale frequency would therefore make a blocking layer look like improved adaptation. Completion accuracy prevents that interpretation.\par
The simulation uses \SimulationCalls{} public scoring calls. It creates no additional tool families and collects no model responses. Its intervention is known by construction. The experiment tests the measurement response to that intervention, not the prevalence of obsolete calls in real deployments. It is a controlled sensitivity analysis, with an assumption-driven response curve rather than a learned-model performance curve.\par
\section{Coverage and task construction}
The authored families cover issue creation, charges, messages, calendar events, support tickets, transactional email, object uploads, keyword searches, text messages, deployments, contact updates, and refunds. These domains provide different argument structures and routing choices. They are not sampled from a complete inventory of deployed APIs. Equal family coverage is a design choice, not an estimate of how often agents use each operation.\par
Each family has a base schema and request pattern. A seeded generator chooses values and a migration variant. It then constructs the current interface and both reference calls. Values include identifiers, messages, categories, routing fields, and numeric quantities. The migration may change representation while preserving the request's intended meaning. The obsolete reference remains a valid description of the old interface, even though it fails the current task.\par
\begin{figure}[t]\centering
\begin{tikzpicture}\begin{axis}[ybar,ymin=0,ymax=15,ylabel={Public tasks},xtick={0,1,2,3,4},xticklabels={Rename,Required,Enum,Endpoint,Type},bar width=28pt]
\addplot[fill=okblue] table[x=x,y=n,col sep=comma]{figures/coverage.generated.csv};
\end{axis}\end{tikzpicture}
\caption{The public set balances the migration categories. Counts describe authored coverage, not the frequency of changes in deployed services.}\label{fig:coverage}
\end{figure}
\paragraph{Field renaming.} A request value moves to a new argument key while its meaning stays fixed. The current property description lets the solver identify the destination. An old call may contain all the right values but place one under an obsolete name. Exact argument matching catches both the missing current field and the unexpected old field. This differs from a semantic misunderstanding of the requested value.
\paragraph{New required information.} A migrated operation requires an additional field. The request explicitly supplies its value. The intended behavior is to add that field while preserving the other arguments. The benchmark does not ask the agent to guess missing information or contact the user. A failure here can reflect omission, extraction failure, or reliance on the old argument set.
\paragraph{Enumeration change.} The request retains old terminology while the current schema states a mapping to new values. Copying the request literally can therefore produce an obsolete call. The task tests whether the agent applies the supplied mapping. It does not require discovering an undocumented relationship between category systems. That restriction makes correctness decidable but simplifies real migrations with ambiguous meaning.
\paragraph{Endpoint split.} A former routing field becomes a choice among operations. The current endpoint descriptions carry the routing information. The intended answer selects the matching operation and omits the removed field. A familiar old operation name is insufficient. Choosing the first listed operation also fails when the request needs another route, and empty arguments remain incomplete regardless of endpoint order.
\paragraph{Type change.} A value must use a new representation, such as a string, a list representation, a duration, or a converted monetary unit. The task checks the resulting value and type, rather than accepting any syntactically valid object. This category is broader than simple type checking. A wrong numeric value can satisfy a schema while still failing the expected-call comparison.
Figure~\ref{fig:coverage} shows the category counts. Uniform coverage supports balanced category inspection, but it does not imply uniform difficulty. Some mutations change a single key; others require a representation conversion. A pooled score hides that distinction. Future learned-model reports should therefore preserve category and family breakdowns, alongside the pooled measure.\par
\section{Failure taxonomy}
The primary outcome separates correct current calls from failures. The diagnostic then separates exact stale reuse from other failures. This hierarchy is intentionally conservative. An answer that nearly reproduces an old call can fail without receiving the exact-stale label. Conversely, an answer can match the old call by coincidence or by misunderstanding the current instructions. The label records output structure, not a causal account of memory.\par
A useful error analysis distinguishes operation errors, field errors, value errors, and missing calls. Operation errors select an obsolete or incorrect endpoint. Field errors omit required arguments or preserve removed keys. Value errors use an obsolete enum, incorrect type, wrong quantity, or another incorrect argument. Missing calls include refusal, empty output, and prose without an extractable call. These descriptions can overlap, so they should not be interpreted as a mutually exclusive learned-model taxonomy without further annotation.\par
The study retains the grader's first-difference explanation for inspection. A first difference is useful for debugging but is not a complete count of all defects. An answer can omit a required field and include an extra field at the same time. Reporting only the first reason would undercount the second defect. The headline statistics therefore use binary correctness and the exact-stale diagnostic, which have explicit definitions.\par
This taxonomy also distinguishes agent error from enforcement behavior. A validator can intercept a bad proposal, but interception does not convert that proposal into a correct one. If the final answer becomes empty, the output taxonomy changes from stale reuse to missing call. That change can be beneficial for containment, but it is not evidence of completed work. The simulated results show why both stages should remain observable.\par
\section{Metrics and scoring}
\subsection{Call extraction and normalization}
The scorer extracts call objects from the answer. It accepts direct name-and-arguments objects and supported function-call wrappers. Arguments may arrive as an object or as a JSON-encoded object. If several extractable calls appear, the last valid call is used. This behavior accommodates common output forms, but it also means the metric does not establish that an earlier harmful proposal was absent.\par
Canonicalization sorts object keys and preserves array order. Normalization treats equivalent supported duration strings alike and removes whitespace around comma-separated list elements. Other value distinctions remain significant. The expected call is compared after the same normalization. A match receives credit; an exact normalized match to the obsolete reference receives no credit and sets the stale diagnostic.\par
These rules define a restricted equivalence relation. They do not prove semantic equivalence for arbitrary APIs. A service might accept an omitted optional argument, another operation sequence, or a value representation outside the accepted forms. Conversely, string normalization might collapse distinctions that matter to a particular real service. The present claim concerns the authored benchmark contract. Extending it requires reviewing accepted alternatives against each domain's semantics.\par
\subsection{Outcome definitions}
Let $C_i$ indicate that answer $i$ matches the current call, and let $S_i$ indicate an exact obsolete-call match. For $N$ tasks, the reported current-call accuracy and stale-call rate are\par
\[ A=\frac{1}{N}\sum_i C_i,\qquad R_{\mathrm{stale}}=\frac{1}{N}\sum_i S_i. \]
The two measures are not complements. An empty answer has neither a current nor a stale match. The same holds for an unrelated call. A low stale rate therefore cannot establish successful adaptation. Accuracy must accompany it. For the simulated validation layer, the report also retains the proposed stale indicator and the blocked indicator before final scoring.\par
A paired contrast compares policies on the same tasks. The visible-versus-stale contrast averages accuracy differences, with resampling grouped by tool family. In the simulated control comparison, each agent uses the same task draws under both layers. Accuracy differences are computed within agent before resampling. These comparisons preserve the experimental pairing instead of treating matched observations as independent samples.\par
\subsection{Intervals and dependence}
All reported empirical rates carry uncertainty intervals. Scripted task rates use Wilson intervals at the stated confidence level. These intervals describe a hypothetical task-sampling model, not randomness in a deterministic policy. The public set is authored, balanced, and dependent within families. The intervals must not be read as a guarantee about arbitrary future APIs. We retain them as a transparent small-sample summary alongside exact counts.\par
For the scripted paired contrast, the bootstrap resamples families and retains their mutation siblings. For simulated rates, the bootstrap resamples whole agents, retaining all tasks within each agent. Each bootstrap uses \BootstrapDraws{} draws and percentile endpoints for a 95\% interval. The simulated intervals consequently describe variation among the seeded simulated agents conditional on this task set. They do not include uncertainty about the simulation assumptions, the chosen schemas, or deployment populations.\par
At a boundary, empirical bootstrap intervals can collapse. If every observed agent has the same outcome, resampling those agents creates no variation. A zero-width interval is then a property of the resampling calculation, not proof of universal certainty. The task-level Wilson intervals in the scripted controls remain important context. No significance threshold is used to turn these diagnostic experiments into a model superiority claim.\par
\section{Results}
\subsection{Scripted measurement checks}
The visible solver achieves \VisibleRate{} current-call accuracy. The oracle has the same task outcomes. The stale reference achieves \FloorRate{} accuracy and \StaleRate{} stale detection. The empty, no-operation, refusal, echo, and first-tool controls each achieve \FloorRate{} accuracy. Table~\ref{tab:baseline} gives both measures for every control, so inactivity cannot be mistaken for adaptation.\par
\begin{table}[t]\centering\small\begin{tabular}{lrr}\toprule Policy & Accuracy (95\% Wilson) & Stale match (95\% Wilson)\\\midrule
visible & 100.0\% [94.0, 100.0] & 0.0\% [0.0, 6.0]\\
oracle & 100.0\% [94.0, 100.0] & 0.0\% [0.0, 6.0]\\
stale & 0.0\% [0.0, 6.0] & 100.0\% [94.0, 100.0]\\
empty & 0.0\% [0.0, 6.0] & 0.0\% [0.0, 6.0]\\
noop & 0.0\% [0.0, 6.0] & 0.0\% [0.0, 6.0]\\
refuse & 0.0\% [0.0, 6.0] & 0.0\% [0.0, 6.0]\\
echo & 0.0\% [0.0, 6.0] & 0.0\% [0.0, 6.0]\\
first & 0.0\% [0.0, 6.0] & 0.0\% [0.0, 6.0]\\
\bottomrule\end{tabular}

\caption{Scripted controls on the public set. Brackets show task-level Wilson intervals. Oracle and stale controls intentionally access reference answers.}\label{tab:baseline}\end{table}
\begin{figure}[t]\centering\begin{tikzpicture}\begin{axis}[ybar,ymin=0,ymax=108,ylabel={Accuracy (percent)},xtick={0,1,2,3,4,5,6,7},xticklabels={Visible,Oracle,Stale,Empty,Noop,Refuse,Echo,First},bar width=17pt]
\addplot[fill=okblue,error bars/.cd,y dir=both,y explicit] table[x=x,y=value,y error plus=plus,y error minus=minus,col sep=comma]{figures/baselines.generated.csv};
\end{axis}\end{tikzpicture}\caption{The scripted solver and oracle separate from the negative controls. Error bars are Wilson intervals over public tasks. This is a measurement check, not a model leaderboard.}\label{fig:baseline}\end{figure}
The first hypothesis holds for the registered controls. The visible solver completes each item using the supplied prompt. This provides evidence against missing information within the authored grammar. It does not establish that the grammar is natural, representative, or difficult for general models. The solver was designed for these templates and was already available before this study.\par
The second hypothesis also holds. The obsolete references fail while receiving the expected diagnostic. The visible-versus-stale accuracy contrast is \PairDifference{} in percentage-point units. Its bootstrap interval collapses because the contrast is identical in every observed family. This is a clear separation between designed controls, but it is not a population estimate of how much migration hurts a trained model.\par
All repeated scripted semantic records agree. The repeat check includes task identity, final answer, score, and diagnostic details. Timing changes are excluded. This result supports deterministic reproduction for the exercised route and environment. It does not establish agreement across every remote interface or future dependency version. The preserved provenance identifies the evaluated sources and dependency lock.\par
\subsection{Simulated control-layer ablation}
At the middle adherence setting, simulated current-call accuracy is \SimAccuracy{} under both layers. Simulated emitted stale frequency is \SimStaleBefore{} with pass-through and \SimStaleAfter{} with validation. The simulated blocking rate is \SimBlocked{}. The simulated paired accuracy change is \SimDifference{} in percentage-point units. Thus, the control-layer completion result is a null.\par
\begin{figure}[t]\centering\begin{tikzpicture}\begin{axis}[ymin=0,ymax=105,xmin=0,xmax=1,xlabel={Simulated adherence parameter $q$},ylabel={Simulated rate (percent)},legend style={at={(.5,-.28)},anchor=north,legend columns=1,font=\footnotesize}]
\addplot[okblue,mark=*,error bars/.cd,y dir=both,y explicit] table[x=q,y=accuracy,y error plus=ap,y error minus=am,col sep=comma]{figures/pass.generated.csv};\addlegendentry{Accuracy, both simulated layers (overlap)}
\addplot[okorange,mark=square*,error bars/.cd,y dir=both,y explicit] table[x=q,y=stale,y error plus=sp,y error minus=sm,col sep=comma]{figures/pass.generated.csv};\addlegendentry{Stale emitted, simulated pass-through}
\addplot[black,dashed,mark=triangle*,error bars/.cd,y dir=both,y explicit] table[x=q,y=stale,y error plus=sp,y error minus=sm,col sep=comma]{figures/validate.generated.csv};\addlegendentry{Stale emitted, simulated validation}
\end{axis}\end{tikzpicture}\caption{Simulated population sweep with agent-cluster bootstrap intervals. Validation suppresses stale outputs but leaves simulated accuracy unchanged. Coincident accuracy curves are drawn once; boundary intervals collapse.}\label{fig:simulation}\end{figure}
Figure~\ref{fig:simulation} shows the full simulated sweep. Increasing the simulated adherence parameter increases current-call accuracy. Validation changes the emitted stale curve but not the accuracy curve. No new correct calls appear because the layer has no repair step. The observed null is therefore consistent with the registered mechanism, rather than evidence that all possible control layers are ineffective.\par
The simulated validator rejects the obsolete proposals represented in this public set. That result does not imply complete stale-call detection for every migration. A numeric unit change can retain the same primitive type, for example. A structurally valid call can also carry the wrong identifier or amount. Such cases require semantic validation or outcome checks. The generated public instances support the observed separation, while the general mechanism has a narrower reach.\par
The simulated experiment also clarifies a reporting risk. If an evaluator replaces proposed calls with blocked outputs and publishes only stale frequency, the control can appear to solve migration adherence. Retaining the pre-control proposal reveals that the simulated policy has not changed. Retaining accuracy reveals that completion has not improved. These observations support a measurement recommendation: report proposals, emitted calls, and completion separately when a control layer intervenes.\par
\section{Roadmap}
The first extension is a genuinely matched migration study. Each request should appear with an unchanged interface and with a migrated interface. A model would answer both under a fixed protocol. The paired difference would then estimate the effect of the migration within the tested task distribution. Order effects, prompt differences, and shared content would need explicit controls. The current positive and obsolete references do not supply this counterfactual by themselves.\par
The second extension is evaluation on learned agents. Such a study should freeze the model identity, prompt, output budget, and retry policy. It should preserve raw answers and specify how native tool-call outputs map into the scoring contract. Public-template scores should be accompanied by fresh task instances and new request phrasings. A held-out seed alone changes values and variants, but does not remove familiarity with the underlying family grammar.\par
The third extension is semantic enforcement. The simulated validator checks structure and supported schema constraints. A richer layer might use unit annotations, reject suspicious values, ask for missing information, or repair a call. Those interventions introduce new outcomes: incorrect repair, unnecessary refusal, extra interaction cost, and delayed completion. Evaluating them requires a protocol that records interventions and final task effects rather than only the last proposed call.\par
The fourth extension is task validation beyond the author. Independent readers should check that prompts determine the expected call and that accepted equivalents cover reasonable answers. Domain experts can assess whether migration patterns reflect actual service changes. These checks would support stronger claims about representativeness. Until then, the benchmark remains an explicit synthetic testbed with a limited set of authored patterns.\par
\section{Caveats and interpretation limits}
\paragraph{Synthetic scope.} The authored families are convenient examples, not a probability sample of deployed interfaces. Their balance helps diagnosis but does not justify prevalence estimates. The request patterns expose information more cleanly than many production tasks. Ambiguous documentation, incomplete schemas, service outages, authorization errors, and conflicting versions are outside the present design.
\paragraph{Cause is not observed.} An exact obsolete call is compatible with several causes. A model might recall an older interface, ignore part of the prompt, misread a mapping, or guess. The benchmark observes an answer, not internal memory. Claims about memorization require additional interventions and evidence. The phrase stale reuse refers only to the exact output diagnostic in this paper.
\paragraph{No learned-model evidence.} The study uses scripted policies and a simulated population. Free hosted inference is not configured, and provider availability is untested. No score is reported for a commercial or open-weight language model. The results cannot establish model difficulty, ranking stability, or expected deployment error rates. The specialized visible solver is a solvability control, not evidence that general models find the tasks easy.
\paragraph{Simulation assumptions.} The simulated population chooses between reference calls. It has no partial adaptation, unrelated errors, uncertainty estimates, or language generation. Its adherence parameter is assigned, not measured. The simulated bootstrap intervals are conditional on that mechanism and the public tasks. They should never be displayed as uncertainty about real agent populations. The simulated null applies to a blocking-only layer, not to repair or retraining.
\paragraph{Scoring limits.} Exact normalized-call comparison is transparent but restrictive. It can reject an alternative call that a real service would accept. It can also ignore earlier proposals because extraction selects the last valid call. The benchmark measures a final answer under its contract, not a complete execution safety property. Future trajectory-based claims must record and score each attempted action.
\paragraph{Registration and independence.} The plan was committed before the study outputs, after inspection of existing code and controls. It was not registered with an external service. The author designed the benchmark and the study. No independent human annotation, expert survey, or external replication is reported. Automated reproduction checks do not substitute for those forms of validation.
\paragraph{Release boundaries.} The public tasks and generator can enter training data. The benchmark supports separately generated task instances, but this paper evaluates the public set only. Reproducing the reported results is different from establishing uncontaminated model performance. Maintenance should preserve old artifacts, document revised scoring rules, and distinguish changes to the task distribution from changes to a model.
\section{Citation, availability, and authors}
Code, data, the study plan, and reproduction instructions are available at \url{https://github.com/pooriaarab/benchmark-schema-mutation}. Cite this paper by its title and author using the accompanying citation metadata. The paper and generated figures use CC-BY-4.0; the benchmark code retains its existing license. The study makes no publication or archival-deposit claim.\par
Pooria Arab is an independent researcher and the benchmark author. AI assistance supported experiment scripting, drafting, and artifact checks. The author remains responsible for the scientific claims and release decisions. There are no human participants or collected personal records in this study. Identifiers in the tasks are synthetic examples. The main evidence concerns the benchmark's measurement behavior, with deployment validity left for future work.\par
\bibliographystyle{plainnat}
\bibliography{references}
\end{document}
